\unit{Exposé 149 — Théorèmes de dualité pour les faisceaux algébriques cohérents}
\sid{149}
\src{169}
\seminar{(Mai 1957)}
\paperhead{THÉORÈMES DE DUALITÉ POUR LES FAISCEAUX ALGÉBRIQUES COHÉRENTS}%
  {par Alexandre GROTHENDIECK}

Les résultats qui suivent, inspirés par le \fquote{théorème de dualité algébrique}
de Serre, ont été trouvés en hiver 1955 et hiver 1956. Ils s'établissent très
simplement à l'aide de résultats assez élémentaires sur la cohomologie des espaces
projectifs [3], et l'utilisation intensive de l'algèbre homologique de
Cartan-Eilenberg, sous la forme [2].

\fsection{\normalfont1}{\normalfont\underline{Les} Ext \underline{de faisceaux de modules} ({[2]}, chap. 3 et 4).}

Soit $\symup{X}$ un espace topologique muni d'un faisceau $\symup{\Osh}$ d'anneaux avec unité
(non nécessairement commutatifs). On considère la catégorie abélienne
$\symup{\Ccat^{\Osh}}$ des faisceaux de $\symup{\Osh}$-modules, appelés aussi $\symup{\Osh}$-Modules.
On sait que tout objet de cette catégorie admet une résolution injective, ce qui
permet d'y définir les foncteurs $\symup{\operatorname{Ext}}$ ayant les propriétés
formelles bien connues. De façon précise, pour éviter des confusions, nous
désignons par $\symup{\operatorname{Hom}_{\Osh}(X;A,B)}$ ou simplement
$\symup{\operatorname{Hom}(X;A,B)}$ le \underline{groupe} abélien des $\symup{\Osh}$-homomorphismes de $\symup{A}$
dans $\symup{B}$, tandis que $\symup{\Homsh_{\Osh}(A,B)}$ désignera le \underline{faisceau} des germes
d'homomorphismes de $\symup{A}$ dans $\symup{B}$ ($\symup{A,B\in\Ccat^{\Osh}}$). On définit pour
$\symup{A\in\Ccat^{\Osh}}$ fixé des foncteurs $\symup{h_A}$ et $\symup{\hsh_A}$ à valeurs respectivement
dans la catégorie $\symup{\Ccat}$ des groupes abéliens et la catégorie
$\symup{\Ccat^{\mathrm Z}}$ des faisceaux abéliens sur $\symup{X}$, par les formules :
\begin{equation}\tag{1.1}
  \symup{h_A(B)=\operatorname{Hom}_{\Osh}(X;A,B)}
  \qquad
  \symup{\hsh_A(B)=\Homsh_{\Osh}(A,B)}
\end{equation}

Les foncteurs $\symup{h_A}$ et $\symup{\hsh_A}$ sont des foncteurs covariants exacts à gauche,
dont on considère les foncteurs dérivés droits, notés respectivement
$\symup{\operatorname{Ext}_{\Osh}^{p}(X;A,B)}$ et $\symup{\Extsh_{\Osh}^{p}(A,B)}$. On a donc
par définition
\begin{equation}\tag{1.2}
\left\{
\begin{aligned}
\symup{\operatorname{Ext}_{\Osh}^{p}(X;A,B)}
  &= \symup{(R^{p}h_A)(B)
   = H^{p}\bigl(\operatorname{Hom}_{\Osh}(X;A,C(B))\bigr)}\\
\symup{\Extsh_{\Osh}^{p}(A,B)}
  &= \symup{(R^{p}\hsh_A)(B)
   = H^{p}\bigl(\Homsh_{\Osh}(A,C(B))\bigr)}
\end{aligned}
\right.
\end{equation}

où le symbole $\symup{R^{p}}$ désigne le passage aux foncteurs dérivés droits, et où
$\symup{C(B)}$ désigne une résolution injective arbitraire de $\symup{B}$ dans $\symup{\Ccat^{\Osh}}$.
Désignons par
$\symup{\Gamma:\Ccat^{\mathrm Z}\longrightarrow\Ccat}$ le foncteur \fquote{sections}. Rappelons que ses foncteurs dérivés droits sont
notés par $\symup{B\longrightarrow H^{p}(X,B)}$ :


\src{170}
\begin{equation}\tag{1.3}
  \symup{H^{p}(X,B)=(R^{p}\Gamma)(B)=H^{p}\bigl(\Gamma(C(B))\bigr)}
\end{equation}
On a évidemment $\symup{h_A=\Gamma\hsh_A}$ ; d'autre part on vérifie que $\symup{\hsh_A}$
transforme objets injectifs en objets acycliques pour $\symup{\Gamma}$. On en conclut de
façon bien connue :

\result{PROPOSITION}{1} \uline{Il existe pour tout} $\symup{\Osh}$-\uline{Module} $\symup{A}$ \uline{un foncteur
spectral cohomologique sur} $\symup{\Ccat^{\Osh}}$, \uline{aboutissant au foncteur gradué}
$\symup{\bigl(\operatorname{Ext}_{\Osh}^{*}(X;A,B)\bigr)}$, \uline{et dont le terme initial est}
\begin{equation}\tag{1.4}
  \symup{E_{2}^{p,q}(A,B)=H^{p}\bigl(X,\Extsh_{\Osh}^{q}(A,B)\bigr)}
\end{equation}
On en déduit des \fquote{\uline{edge-homomorphisms}} et une suite exacte à cinq termes
que nous n'écrirons pas.

\result{COROLLAIRE}{1} \uline{Si} $\symup{A}$ \uline{est localement isomorphe à} $\symup{\Osh^{n}}$, \uline{alors on a
des isomorphismes canoniques}
\begin{equation}\tag{1.5}
  \symup{\operatorname{Ext}_{\Osh}^{p}(X;A,B)
  \overset{\sim}{\longleftarrow}
  H^{p}\bigl(X,\Homsh_{\Osh}(A,B)\bigr)}
\end{equation}
(donnés par des edge-homomorphisms de la suite spectrale). \uline{En particulier on a
un isomorphisme canonique}
\begin{equation}\tag{1.6}
  \symup{\operatorname{Ext}_{\Osh}^{p}(X;\Osh,B)=H^{p}(X,B)}
\end{equation}

Pour utiliser ces résultats, il faut savoir expliciter les
$\symup{\Extsh_{\Osh}^{p}(A,B)}$. Ce sont des foncteurs qui se calculent localement,
i.e. si $\symup{U}$ est un ouvert de $\symup{X}$, on a
\[
  \symup{\Extsh_{\Osh}^{p}(A,B)|U
  =\Extsh_{\Osh|U}^{p}(A|U,B|U)}
\]
comme il résulte du fait que la restriction à $\symup{U}$ d'un $\symup{\Osh}$-Module injectif
est un $\symup{(\Osh|U)}$-Module injectif. De plus, on a des homomorphismes fonctoriels,
pour $\symup{x\in X}$ fixé
\begin{equation}\tag{1.7}
  \symup{\Homsh_{\Osh}(A,B)_x
  \longrightarrow
  \operatorname{Hom}_{\Osh_x}(A_x,B_x)}
\end{equation}
qui se prolongent de façon unique en un homomorphisme de foncteurs cohomologiques
en $\symup{B}$ :
\begin{equation}\tag{1.8}
  \symup{\Extsh_{\Osh}^{p}(A,B)_x
  \longrightarrow
  \operatorname{Ext}_{\Osh_x}^{p}(A_x,B_x)}
\end{equation}

\result{PROPOSITION}{2} \uline{Si} $\symup{A}$ \uline{est, dans un voisinage de} $\symup{x}$, \uline{isomorphe au
conoyau d'un homomorphisme} $\symup{\Osh^{m}\longrightarrow\Osh^{n}}$, \uline{alors} (1.7) \uline{est un isomorphisme pour tout} $\symup{p}$. \uline{Il en est ainsi en
particulier si} $\symup{A}$ \uline{est un} $\symup{\Osh}$-\uline{Module cohérent} [3].

\src{171}
\result{PROPOSITION}{3} \uline{Soit} $\symup{L_{*}=(L_i)}$ \uline{une résolution gauche du}
$\symup{\Osh}$-\uline{Module} $\symup{A}$ \uline{par des} $\symup{\Osh}$-\uline{Modules dont chacun est isomorphe localement à
un} $\symup{\Osh^{n}}$. \uline{Alors} $\symup{\Extsh_{\Osh}(A,B)}$ \uline{s'identifie à}
$\symup{H^{*}\bigl(\Homsh_{\Osh}(L_{*},B)\bigr)}$, \uline{et}
$\symup{\operatorname{Ext}_{\Osh}(X;A,B)}$ \uline{s'identifie à l'hypercohomologie de} $\symup{X}$
\uline{par rapport au complexe} $\symup{\Homsh_{\Osh}(L_{*},B)}$.

La démonstration est standard, on considère le bicomplexe
$\symup{\Homsh_{\Osh}(L_{*},C(B))}$ où $\symup{C(B)}$ est une résolution injective de $\symup{B}$, et les
homomorphismes naturels de $\symup{\Homsh_{\Osh}(L_{*},B)}$ et de
$\symup{\Homsh_{\Osh}(A,C(B))}$ dans ce dernier.

Pour finir, indiquons que les deux foncteurs $\symup{\operatorname{Ext}}$ introduits
dans (1.2) sont non seulement des foncteurs cohomologiques en $\symup{B}$, mais des
\uline{bifoncteurs cohomologiques}, covariants en $\symup{B}$ et contrevariants en $\symup{A}$.

\fsection{\normalfont2}{\normalfont\uline{Loi de composition dans les} Ext.}

Les résultats de ce numéro sont dus indépendamment à CARTIER et à YONEDA ; voir
un exposé de CARTIER [1] pour des détails. Soit $\symup{\underline{C}}$ une catégorie abélienne.
Soient $\symup{K}$ et $\symup{L}$ deux objets gradués dans $\symup{\underline{C}}$. On désigne par
$\symup{\operatorname{Hom}(K,L)}$ le groupe abélien gradué dont la composante de degré
$\symup{n}$ est formée des homomorphismes homogènes de degré $\symup{n}$ de $\symup{K}$ dans $\symup{L}$ (i.e.
des systèmes $\symup{(u_i)}$ d'homomorphismes $\symup{K^{i}\longrightarrow L^{i+n}}$).
Si $\symup{K}$ et $\symup{L}$ sont des complexes (à opérateurs différentiels de degré $\symup{+1}$ pour
fixer les idées), on introduit dans $\symup{\operatorname{Hom}(K,L)}$ l'opérateur
différentiel donné par
\begin{equation}\tag{2.1}
  \symup{\delta(u)=du+(-1)^{n+1}ud,\qquad\text{où }n=\deg(u)}
\end{equation}
qui en fait un complexe à opérateur différentiel de degré $\symup{+1}$. Les cycles de
degré $\symup{n}$ sont les opérateurs de degré $\symup{n}$ qui commutent à $\symup{d}$ (en tant
qu'opérateurs homogènes).\footnote{Le texte imprimé porte « anticommutent à $\symup{u}$ ». Par (2.1), un cycle vérifie $\symup{du=(-1)^nud}$ : il commute donc à $\symup{d}$ au sens gradué.} On peut alors considérer
$\symup{H^{*}\bigl(\operatorname{Hom}(K,L)\bigr)}$, c'est un invariant des types
d'homotopie de $\symup{K}$ et de $\symup{L}$, qu'on pourra noter $\symup{H^{*}(K,L)}$.
Si on a un troisième complexe $\symup{M}$, alors la composition des homomorphismes
définit un accouplement
$\symup{\operatorname{Hom}(K,L)\times\operatorname{Hom}(L,M)
\longrightarrow\operatorname{Hom}(K,M)}$ compatible avec les opérateurs différentiels, d'où par passage à la cohomologie
des accouplements
\begin{equation}\tag{2.2}
  \symup{H^{*}(K,L)\times H^{*}(L,M)\longrightarrow H^{*}(K,M)}
\end{equation}
notés $\symup{(u,v)\longrightarrow vu}$. Ces accouplements satisfont une propriété évidente
d'associativité. En particulier, $\symup{H^{*}(K,K)}$ est un anneau gradué associatif
avec unité, et $\symup{H^{*}(K,L)}$ (resp. $\symup{H^{*}(L,K)}$) est un module gradué à droite
(resp. à gauche) sur cet anneau, etc. En dimension $\symup{0}$, (2.2) se réduit à la
composition des homomorphismes permis de complexes. Enfin, une suite exacte de
complexes
\src{172}
$\symup{0\longrightarrow K'\longrightarrow K\longrightarrow K''\longrightarrow 0}$, telle que $\symup{K'^{i}}$ s'identifie à un facteur direct de $\symup{K^{i}}$ pour tout $\symup{i}$,
donne naissance à une suite exacte pour les complexes de groupes
$\symup{\operatorname{Hom}(K'',L)}$, etc., d'où un opérateur cobord
$\symup{H^{i}(K',L)\longrightarrow H^{i+1}(K'',L)}$.
On définit de même les opérateurs cobord relatifs à une suite exacte en $\symup{L}$. Les
accouplements (2.2) sont compatibles, au sens usuel, avec ces opérateurs cobord.

Supposons maintenant que $\symup{C}$ soit une catégorie telle que tout élément $\symup{A}$
de $\symup{C}$ admette une résolution injective $\symup{C(A)}$. On constate alors, utilisant
une des nombreuses variantes du théorème du bicomplexe, que
\[
  \symup{H^{*}\bigl(C(A),C(B)\bigr)
  =H^{*}\bigl(\operatorname{Hom}(C(A),C(B))\bigr)}
\]
est canoniquement isomorphe à
\[
  \symup{H^{*}\bigl(\operatorname{Hom}(A,C(B))\bigr)
  =\operatorname{Ext}^{*}(A,B).}
\]
Les opérateurs cobord envisagés plus haut donnent les homomorphismes cobord des
$\symup{\operatorname{Ext}}$. De plus, les accouplements (2.2) donnent ici des
accouplements associatifs, compatibles avec les homomorphismes cobord :
\begin{equation}\tag{2.3}
  \symup{\operatorname{Ext}^{*}(A,B)\times\operatorname{Ext}^{*}(B,C)
  \longrightarrow\operatorname{Ext}^{*}(A,C)}
\end{equation}
En particulier, $\symup{\operatorname{Ext}^{*}(A,A)}$ est un anneau gradué associatif
avec unité, etc. (On montre de façon analogue que les foncteurs
$\symup{\operatorname{Ext}}$ opèrent dans les foncteurs dérivés d'un foncteur
quelconque ; nous ne nous servirons pas ici de ce fait.)

Dans le cas où la catégorie envisagée est la catégorie $\symup{\Ccat^{\Osh}}$ des
$\symup{\Osh}$-Modules sur $\symup{X}$, on obtient donc des accouplements
\begin{equation}\tag{2.4}
  \symup{\operatorname{Ext}_{\Osh}^{p}(X;A,B)
  \times\operatorname{Ext}_{\Osh}^{q}(X;B,C)
  \longrightarrow\operatorname{Ext}_{\Osh}^{p+q}(X;A,C)}
\end{equation}
qui peuvent s'expliciter comme il a été dit. D'ailleurs, les mêmes
développements, mais où on remplace la catégorie des groupes abéliens par la
catégorie des faisceaux abéliens sur $\symup{X}$, et les foncteurs $\symup{\operatorname{Hom}}$
par les foncteurs $\symup{\Homsh}$, définissent aussi des accouplements, ayant les mêmes
propriétés formelles, et de \fquote{nature locale} cette fois-ci :
\begin{equation}\tag{2.5}
  \symup{\Extsh_{\Osh}^{p}(A,B)\times\Extsh_{\Osh}^{q}(B,C)
  \longrightarrow\Extsh_{\Osh}^{p+q}(A,C)}
\end{equation}
Ces derniers se précisent si on remarque que les homomorphismes (1.8) sont
compatibles avec les accouplements entre les $\symup{\operatorname{Ext}}$.

\src{173}
Rappelons enfin qu'on a aussi une structure multiplicative entre foncteurs
$\symup{H^{p}(X,A)}$ (cup-produit). On constate alors que les suites spectrales de la
proposition 1 sont compatibles avec les structures multiplicatives, de façon
plus précise, on a un accouplement de la suite spectrale $\symup{E(A,B)}$ et de la suite
spectrale $\symup{E(B,C)}$ vers la suite spectrale $\symup{E(A,C)}$, qui, pour l'aboutissement
se réduit à l'accouplement entre les $\symup{\operatorname{Ext}}$ globaux, et pour le
terme initial à l'accouplement déduit, dans le deuxième membre de (1.4), du
cup-produit et des accouplements des $\symup{\operatorname{Ext}}$ locaux. Il en résulte
en particulier que les \fquote{edge-homomorphisms}
\begin{equation}\tag{2.6}
  \symup{\operatorname{Ext}_{\Osh}^{n}(X;A,B)
  \longrightarrow H^{0}\bigl(X;\Extsh_{\Osh}^{n}(A,B)\bigr)}
\end{equation}
\begin{equation}\tag{2.7}
  \symup{H^{n}\bigl(X,\Homsh_{\Osh}(A,B)\bigr)
  \longrightarrow\operatorname{Ext}_{\Osh}^{n}(X;A,B)}
\end{equation}
sont compatibles avec les structures multiplicatives. Si on se borne donc aux
faisceaux localement isomorphes à un $\symup{\Osh^{m}}$, cela explicite complètement la
composition entre $\symup{\operatorname{Ext}}$ globaux à l'aide du cup-produit, compte
tenu de l'isomorphisme (1.5).

\fsection{\normalfont3}{\normalfont\uline{Résultats de cohomologie locale}.}

Soit $\symup{A}$ un anneau commutatif avec unité muni d'un idéal $\symup{J}$. Nous allons
définir, pour un $\symup{A}$-module $\symup{M}$ variable, des homomorphismes fonctoriels
\begin{equation}\tag{3.1}
\left\{
\begin{aligned}\symup{\operatorname{Ext}_{A}^{p}(A/J,M)}&\symup{\longrightarrow
  \operatorname{Hom}_{A}\!\left(\bigwedge^{p}J/J^{2},M\otimes A/J\right)}\\\symup{\operatorname{Tor}_{p}^{A}(A/J,M)}&\symup{\longleftarrow
  \left(\bigwedge^{p}J/J^{2}\right)
  \otimes\operatorname{Hom}_{A}(A/J,M)}\end{aligned}
\right.
\end{equation}
où les produits tensoriels et extérieurs sont pris sur l'anneau $\symup{A}$ ; noter
d'ailleurs que $\symup{J/J^{2}}$ est en fait un $\symup{A/J}$-module, et que ses puissances
extérieures en tant que $\symup{A}$-module ou en tant que $\symup{A/J}$-module sont les mêmes.
La définition des homomorphismes (3.1) revient à la définition, pour tout
système $\symup{\underline{x}=(x_1,\ldots,x_p)}$ de points de $\symup{J}$ d'homomorphismes
$\symup{\varphi_{\underline{x}}}$ :
\begin{equation}\tag{3.2}
\left\{
\begin{aligned}\symup{\varphi_{\underline{x}}:
  \operatorname{Ext}_{A}^{p}(A/J,M)}&\symup{\longrightarrow M\otimes A/J}\\\symup{\varphi_{\underline{x}}:
  \operatorname{Hom}_{A}(A/J,M)}&\symup{\longrightarrow
  \operatorname{Tor}_{p}^{A}(A/J,M)}\end{aligned}
\right.
\end{equation}
de telle façon que les conditions suivantes soient satisfaites :

\src{174}
\begin{enumerate}
\item[\textup{i.}] $\symup{\varphi_{x_1,\ldots,x_p}}$ dépend de façon
$\symup{A}$-multilinéaire alternée du système des $\symup{x_i\in J}$
\item[\textup{ii.}] $\symup{\varphi_{x_1,\ldots,x_p}}$ est nul quand un des $\symup{x_i}$ est
dans $\symup{J^{2}}$.
\end{enumerate}
En fait, ii. résulte de i., puisque $\symup{a\varphi_{\underline{x}}=0}$ pour $\symup{a\in J}$,
comme on voit en remarquant que tous les modules dans (3.2) sont annulés par
$\symup{J}$.

Pour définir $\symup{\varphi_{\underline{x}}}$, considérons le complexe
$\symup{K_{\underline{x}}}$ dont le $\symup{A}$-module sous-jacent est $\symup{\bigwedge A^{p}}$, et
dont l'opérateur différentiel est le produit intérieur $\symup{i_{\underline{x}}}$ par
$\symup{\underline{x}}$ considéré comme forme linéaire sur $\symup{A^{p}}$ de composantes
$\symup{x_1,\ldots,x_p}$. L'opérateur différentiel est de degré $\symup{-1}$, les degrés sont
positifs, et la cohomologie de ce complexe en dimension $\symup{0}$ est
$\symup{A/(x_1A+\cdots+x_pA)}$. Comme les $\symup{x_i}$ sont dans $\symup{J}$, on en déduit une
augmentation $\symup{K_{\underline{x},0}\longrightarrow A/J}$. Ainsi
$\symup{K_{\underline{x}}}$ apparaît comme un complexe \uline{libre} augmenté, à module
d'augmentation $\symup{A/J}$. On en déduit, de façon connue, des homomorphismes
\[
\symup{\operatorname{Ext}_{A}^{*}\bigl(H_0(K_{\underline{x}}),M\bigr)
\longrightarrow
H^{*}\bigl(\operatorname{Hom}_{A}(K_{\underline{x}},M)\bigr)\quad\text{et}\quad
\operatorname{Tor}_{*}^{A}\bigl(H_0(K_{\underline{x}}),M\bigr)
\longleftarrow
H_{*}\bigl(K_{\underline{x}}\otimes M\bigr)}
\]
d'où, en composant avec les homomorphismes sur les $\symup{\operatorname{Ext}}$ et les
$\symup{\operatorname{Tor}}$ déduits de l'homomorphisme d'augmentation
$\symup{H_0(K_{\underline{x}})\longrightarrow A/J}$, des homomorphismes
\begin{equation}\tag{3.3}
\left\{
\begin{aligned}\symup{\psi_{\underline{x}}:
\operatorname{Ext}_{A}^{*}(A/J,M)}&\symup{\longrightarrow H^{*}\bigl(\operatorname{Hom}_{A}(K_{\underline{x}},M)\bigr)}\\\symup{\psi_{\underline{x}}:
\operatorname{Tor}_{*}^{A}(A/J,M)}&\symup{\longleftarrow H_{*}\bigl(K_{\underline{x}}\otimes M\bigr)}\end{aligned}
\right.
\end{equation}
Or on constate aussitôt qu'en dimension maxima $\symup{p}$, la cohomologie des seconds
membres est $\symup{M/(x_1M+\cdots+x_pM)}$ (resp. est l'ensemble des éléments de $\symup{M}$
annulés par chacun des $\symup{x_i}$). Comme les $\symup{x_i}$ sont dans $\symup{J}$, on en déduit des
homomorphismes
\begin{equation}\tag{3.4}
\left\{
\begin{aligned}\symup{H^{p}\bigl(\operatorname{Hom}_{A}(K_{\underline{x}},M)\bigr)}&\symup{\longrightarrow M\otimes A/J}\\\symup{H_{p}\bigl(K_{\underline{x}}\otimes M\bigr)}&\symup{\longleftarrow\operatorname{Hom}_{A}(A/J,M)}\end{aligned}
\right.
\end{equation}
Composant les homomorphismes (3.3) et (3.4), on obtient les homomorphismes (3.2)
que nous voulions définir. La vérification de i. est fastidieuse, mais ne
présente pas de difficultés.

\src{175}
\result{PROPOSITION}{4} \uline{Soit} $\symup{A}$ \uline{un anneau commutatif avec unité, soit}
$\symup{(x_1,\ldots,x_p)}$ \uline{une suite d'éléments de} $\symup{A}$ \uline{telle que pour} $\symup{1\leq i\leq p}$,
\uline{l'image de} $\symup{x_i}$ \uline{dans le quotient de} $\symup{A}$ \uline{par l'idéal engendré par}
$\symup{(x_1,\ldots,x_{i-1})}$ \uline{ne soit pas diviseur de} $\symup{0}$. \uline{Soit} $\symup{J}$ \uline{l'idéal engendré
par les} $\symup{x_i}$. \uline{Alors} $\symup{J/J^{2}}$ \uline{est un} $\symup{(A/J)}$\uline{-module libre ayant pour base les
images canoniques des} $\symup{x_i}$, \uline{le complexe} $\symup{K_{\underline{x}}}$ \uline{est une résolution
libre de} $\symup{A/J}$, \uline{et pour tout} $\symup{A}$\uline{-module} $\symup{M}$, \uline{les homomorphismes} (3.1) \uline{en
dimension} $\symup{p}$ \uline{sont bijectifs. Il en est de même des homomorphismes analogues
définis pour des degrés} $\symup{i}$ \uline{quelconques pourvu que} $\symup{J.M=0}$.

(Le point essentiel, dont tous les autres résultent, est l'acyclicité de
$\symup{K_{\underline{x}}}$, qui est un fait bien connu sous les conditions indiquées).

\result{COROLLAIRE}{1} $\symup{A}$ \uline{et} $\symup{J}$ \uline{étant comme ci-dessus, supposons de plus que}
$\symup{A}$ \uline{soit une algèbre affine régulière de} $\symup{\dim n}$ \uline{sur un corps parfait} $\symup{\underline{k}}$, \uline{et
que} $\symup{A/J}$ \uline{soit une algèbre affine régulière. Désignons par} $\symup{\Omega^{1}(A)}$,
$\symup{\Omega^{1}(A/J)}$ \uline{les modules de différentielles de Kähler. Alors on a un
isomorphisme canonique compatible avec la localisation},
\begin{equation}\tag{3.5}
  \symup{\operatorname{Ext}_{A}^{p}\bigl(\Omega^{n-p}(A/J),\Omega^{n}(A)\bigr)=A/J}
\end{equation}
En effet, $\symup{\Omega^{n-p}(A/J)}$ est un $\symup{(A/J)}$-module projectif de rang $\symup{1}$, de même
$\symup{\Omega^{n}(A)}$ est un $\symup{A}$-module projectif de rang $\symup{1}$,\footnote{Le texte imprimé dit « libre » pour les deux modules et « rang $\symup{n}$ » pour $\symup{\Omega^{n}(A)}$. Les modules de formes de degré maximal sont projectifs de rang $\symup{1}$ ; leur liberté globale ne découle pas des hypothèses. L'identification ci-dessous résulte de la suite conormale en prenant les déterminants.}
donc le premier membre est égal à
\[
  \symup{\operatorname{Ext}_{A}^{p}(A/J,A)
  \otimes\Omega^{n-p}(A/J)'\otimes\Omega^{n}(A)}
\]
(où le symbole \fquote{$\symup{\prime}$} désigne le $\symup{(A/J)}$-module dual). Le produit tensoriel des
deux derniers facteurs s'identifie à $\symup{\bigwedge^{p}(J/J^{2})}$, donc le tout
s'identifie à
$\symup{\operatorname{Ext}_{A}^{p}\!\left(A/J,\bigwedge^{p}(J/J^{2})\right)}$,
donc, en vertu de la proposition, à
\[
  \symup{\operatorname{Hom}_{A}\!\left(\bigwedge^{p}J/J^{2},
  \bigwedge^{p}J/J^{2}\right),}
\]
i.e. à $\symup{A/J}$. En particulier, il y a dans
$\symup{\operatorname{Ext}_{A}^{p}\bigl(\Omega^{n-p}(A/J),\Omega^{n}(A)\bigr)}$ un
élément privilégié, correspondant à l'élément unité de $\symup{A/J}$, appelé
\uline{classe fondamentale} de l'idéal $\symup{J}$ dans $\symup{A}$. (Elle peut en fait se
définir sous des conditions sensiblement plus larges). On peut mettre le
corollaire 1 sous une forme plus géométrique et plus globale :

\result{COROLLAIRE}{2} \uline{Soient} $\symup{X}$ \uline{une variété non singulière sur un corps
algébriquement clos} $\symup{\underline{k}}$, $\symup{Y}$ \uline{une sous-variété fermée non singulière de} $\symup{X}$,
$\symup{\Osh_X}$ \uline{le faisceau des anneaux locaux de} $\symup{X}$, $\symup{\Osh_Y}$ \uline{le faisceau des anneaux
locaux de} $\symup{Y}$ \uline{considéré comme un faisceau quotient de} $\symup{\Osh_X}$. \uline{Soient} $\symup{n}$ \uline{la
dimension de} $\symup{X}$, $\symup{n-p}$ \uline{celle de} $\symup{Y}$.

\src{176}
\uline{Soit} $\symup{\underline{\Omega}_X}$ (resp. $\symup{\underline{\Omega}_Y}$) \uline{le faisceau des germes de formes
différentielles régulières sur} $\symup{X}$ (resp. $\symup{Y}$). \uline{Alors on a des isomorphismes
canoniques} :
\begin{equation}\tag{3.6}
  \symup{\Extsh_{\Osh_X}^{p}\bigl(\underline{\Omega}_Y^{n-p},\underline{\Omega}_X^{n}\bigr)=\Osh_Y}
\end{equation}
\uline{ou encore}
\begin{equation}\tag{3.6 bis}
  \symup{\Extsh_{\Osh_X}^{p}\bigl(\Osh_Y,\underline{\Omega}_X^{n}\bigr)=\underline{\Omega}_Y^{n-p}}
\end{equation}
La formule (3.6 bis) peut servir de \uline{définition} de $\symup{\underline{\Omega}_Y^{n-p}}$ pour $\symup{Y}$
variété singulière. De façon précise :

\result{PROPOSITION}{5} \uline{Soit} $\symup{Y}$ \uline{un sous-ensemble algébrique de dimension}
$\symup{q=n-p}$ \uline{dans une variété algébrique non singulière} $\symup{X}$ \uline{de dimension} $\symup{n}$. \uline{Soient}
$\symup{F}$ \uline{un faisceau algébrique cohérent sur} $\symup{X}$ \uline{de support contenu dans} $\symup{Y}$, \uline{et} $\symup{L}$
\uline{un faisceau algébrique localement libre sur} $\symup{X}$. \uline{Alors les faisceaux}
$\symup{\Extsh_{\Osh_X}^{i}(F,L)}$ \uline{sont nuls pour} $\symup{i<p}$ ; \uline{tandis que pour} $\symup{i=p}$, \uline{on a un
isomorphisme canonique}
\begin{equation}\tag{3.7}
  \symup{\Extsh_{\Osh_X}^{p}(F,L)
  =\Homsh_{\Osh_X}\!\left(F,
  \Extsh^{p}(\Osh_X/J,L)\right)}
\end{equation}
\uline{où} $\symup{J}$ \uline{désigne un faisceau d'idéaux arbitraire sur} $\symup{X}$ \uline{annulant} $\symup{F}$ \uline{et dont
l'ensemble des zéros soit} $\symup{Y}$. \uline{En particulier, si} $\symup{F}$ \uline{est un faisceau algébrique
cohérent sur} $\symup{Y}$, \uline{on a}
\begin{equation}\tag{3.7 bis}
  \symup{\Extsh_{\Osh_X}^{p}(F,L)
  =\Homsh_{\Osh_Y}\!\left(F,\Extsh^{p}(\Osh_Y,L)\right).}
\end{equation}
\uline{Enfin}, $\symup{F}$ \uline{étant toujours un faisceau algébrique cohérent sur} $\symup{Y}$, \uline{les faisceaux}
$\symup{E^{i}(F)=\Extsh_{\Osh_X}^{p+i}(F,\underline{\Omega}_X^{n})}$ \uline{ne dépendent pas de la façon d'immerger l'espace algébrique} $\symup{Y}$ \uline{dans une variété
non singulière} $\symup{X}$.

La question étant locale, on peut supposer $\symup{X}$ affine et $\symup{L=\Osh_X}$. Cela ramène
à une question d'algèbre commutative, et même d'algèbre locale : Si $\symup{A}$ est une
localité régulière, $\symup{M}$ un $\symup{A}$-module dont le support est de dimension
$\symup{\leq q=n-p}$ prouver que $\symup{\operatorname{Ext}_{A}^{i}(M,A)=0}$ pour $\symup{i<p}$, et
$\symup{\operatorname{Ext}_{A}^{p}(M,A)
=\operatorname{Hom}_{A}\!\left(M,\operatorname{Ext}^{p}(A/J,A)\right)}$,
où $\symup{J}$ est un idéal quelconque de \fquote{dimension} $\symup{\leq q}$ annulant $\symup{M}$. Pour le
premier point, on procède par récurrence sur $\symup{q}$ : un dévissage immédiat ramène
au cas où $\symup{M}$ est de la forme $\symup{A/J}$, puis, en remplaçant $\symup{J}$ par un idéal plus
petit et utilisant l'hypothèse de récurrence, et la suite exacte des
$\symup{\operatorname{Ext}}$, au cas où $\symup{J}$ est engendré par un \fquote{système de paramètres}
comme dans la proposition 4 et où le résultat est immédiat. Le résultat
précédent implique que si $\symup{J}$ est un idéal

\src{177}
fixé de \fquote{dimension} $\symup{\leq q}$, alors le foncteur contravariant
$\symup{E(M)=\operatorname{Ext}_{A}^{p}(M,A)}$ sur la catégorie des $\symup{A/J}$-modules est
exact à gauche ; de plus il transforme sommes directes en produits directs,
d'où résulte facilement que
$\symup{E(M)=\operatorname{Hom}_{A}\bigl(M,E(A/J)\bigr)}$.
\footnote{Le texte imprimé porte $\symup{E(A)}$ ; dans la catégorie considérée, l'objet
représentant est $\symup{A/J}$, et l'égalité requise pour (3.7) est $\symup{E(A/J)}$.} Enfin, la
dernière assertion de la proposition 5 est plus subtile, et résulte d'une
caractérisation intrinsèque des $\symup{E^{i}(F)}$ à l'aide d'un théorème de dualité
locale qui ne peut être donné ici.

\result{COROLLAIRE}{} \uline{Désignons par} $\symup{\underline{\omega}_Y^{q}}$ \uline{le faisceau}
$\symup{\Extsh_{\Osh_X}^{p}(\Osh_Y,\underline{\Omega}_X^{n})}$. \uline{Alors on a un isomorphisme
fonctoriel pour les faisceaux algébriques cohérents} $\symup{F}$ \uline{sur} $\symup{Y}$ :
\begin{equation}\tag{3.8}
  \symup{\Extsh_{\Osh_X}^{p}(F,\underline{\Omega}_X^{n})
  =\Homsh_{\Osh_X}(F,\underline{\omega}_Y^{q})}
\end{equation}

\fsection{\normalfont4}{\normalfont\uline{Classe de cohomologie associée à une sous-variété}.}

Dans toute la suite, $\symup{X}$ désigne un ensemble algébrique de dimension $\symup{n}$, défini
sur un corps $\symup{\underline{k}}$, que nous supposerons algébriquement clos pour simplifier. Sauf
au n\textsuperscript{o} 6, $\symup{X}$ est supposée non singulière. On désigne par
$\symup{\Osh_X}$ le faisceau des anneaux locaux sur $\symup{X}$, par
$\symup{\underline{\Omega}_X^{*}=\bigoplus\limits_p\underline{\Omega}_X^{p}}$ le faisceau des germes de formes différentielles sur $\symup{X}$.\footnote{Le texte
imprimé emploie un signe de réunion ; la structure graduée est la somme directe
des faisceaux $\symup{\underline{\Omega}_X^p}$.} Si $\symup{Y}$ est une partie fermée de $\symup{X}$, on identifie
les faisceaux algébriques cohérents sur $\symup{Y}$ à des faisceaux cohérents sur $\symup{X}$
nuls hors de $\symup{Y}$ ; il en est en particulier ainsi de $\symup{\Osh_Y}$ et $\symup{\underline{\Omega}_Y}$.

\result{LEMME}{1} \uline{Soient} $\symup{F}$ \uline{un faisceau algébrique cohérent sur} $\symup{X}$ \uline{dont le
support est de dimension} $\symup{\leq n-p}$, $\symup{L}$ \uline{un faisceau algébrique cohérent
localement libre sur} $\symup{X}$. \uline{Alors} $\symup{\operatorname{Ext}_{\Osh_X}^{i}(X;F,L)}$ \uline{est nul
pour} $\symup{i<p}$, \uline{et on a un isomorphisme canonique}
\begin{equation}\tag{4.1}
  \symup{\operatorname{Ext}_{\Osh}^{p}(X;F,L)
  =H^{0}\bigl(X,\Extsh_{\Osh}^{p}(F,L)\bigr)}
\end{equation}
\uline{Si} $\symup{F}$ \uline{est un faisceau algébrique cohérent sur une partie fermée} $\symup{Y}$ \uline{de} $\symup{X}$ \uline{de
dimension} $\symup{\leq n-p}$, \uline{on a un isomorphisme canonique}
\begin{equation}\tag{4.1 bis}
  \symup{\Extsh_{\Osh_X}^{p}(F,L)
  =\Homsh_{\Osh_X}\!\left(
    F\otimes L'\otimes\underline{\Omega}_X^{n},\underline{\omega}_Y^{n-p}
  \right)}
\end{equation}
\uline{où} $\symup{\underline{\omega}_Y^{n-p}}$ \uline{est le faisceau sur} $\symup{Y}$ \uline{défini dans le corollaire à la
proposition 5, (qui s'identifie à} $\symup{\underline{\Omega}_Y^{n-p}}$ \uline{si} $\symup{Y}$ \uline{est non singulière}).

La formule (4.1) est une conséquence immédiate de la suite spectrale de la
proposition 1, et de la proposition 5 ; en vertu de la formule (3.8) on peut
écrire

\src{178}
\[
\symup{\Extsh_{\Osh_X}^{p}(F,L)
=L\otimes(\underline{\Omega}_X^{n})'\otimes
  \Extsh_{\Osh_X}^{p}(F,\underline{\Omega}_X^{n})
=L\otimes(\underline{\Omega}_X^{n})'\otimes
  \Homsh_{\Osh_X}(F,\underline{\omega}_Y^{q})
=\Homsh_{\Osh_X}\!\left(
  F\otimes L'\otimes\underline{\Omega}_X^{n},\underline{\omega}_Y^{q}\right)}
\]
où $\symup{q=n-p}$, d'où la formule (4.1 bis).

Faisons en particulier $\symup{F=\Osh_Y}$, $\symup{L=\underline{\Omega}_X^{p}}$, on trouve, compte tenu
que $\symup{\underline{\Omega}_X^{n}\otimes(\underline{\Omega}_X^{p})'=\underline{\Omega}_X^{n-p}}$, un isomorphisme
canonique
\begin{equation}\tag{4.2}
  \symup{\operatorname{Ext}_{\Osh_X}^{p}(X;\Osh_Y,\underline{\Omega}_X^{p})
  =\operatorname{Hom}_{\Osh_X}
  (X;\underline{\Omega}_X^{n-p},\underline{\omega}_Y^{n-p})}
\end{equation}
Supposons maintenant $\symup{Y}$ \uline{non singulière} pour simplifier, donc
$\symup{\underline{\omega}_Y^{n-p}=\underline{\Omega}_Y^{n-p}}$. On a un homomorphisme naturel de
$\symup{\underline{\Omega}_X^{n-p}}$ dans $\symup{\underline{\Omega}_Y^{n-p}}$, d'où une section canonique $\symup{s_Y}$ du
faisceau $\symup{\Extsh_{\Osh_X}^{p}(\Osh_Y,\underline{\Omega}_X^{p})}$, que nous appellerons, si
toutes les composantes de $\symup{Y}$ sont de dimension $\symup{n-p}$, \uline{section
fondamentale} du faisceau $\symup{\Extsh_{\Osh_X}^{p}(\Osh_Y,\underline{\Omega}_X^{p})}$. Cette
section définit en vertu de (4.1) un élément de
$\symup{\operatorname{Ext}_{\Osh_X}^{p}(X;\Osh_Y,\underline{\Omega}_X^{p})}$. Or l'homomorphisme
naturel $\symup{\Osh_X\longrightarrow\Osh_Y}$ définit un homomorphisme
\[
  \symup{\operatorname{Ext}_{\Osh_X}^{p}(X;\Osh_Y,\underline{\Omega}_X^{p})
  \longrightarrow
  \operatorname{Ext}_{\Osh_X}^{p}(X;\Osh_X,\underline{\Omega}_X^{p})
  =H^{p}(X,\underline{\Omega}_X^{p}).}
\]
On obtient ainsi un élément de $\symup{H^{p}(X,\underline{\Omega}_X^{p})}$, noté $\symup{P_X(Y)}$ et appelé
\uline{classe de cohomologie de} $\symup{Y}$ \uline{dans} $\symup{X}$. Elle est donc déduite de la section
$\symup{s_Y}$ de $\symup{\Extsh_{\Osh_X}^{p}(\Osh_Y,\underline{\Omega}_X^{p})}$ par le diagramme
d'homomorphismes suivants :
\begin{equation}\tag{4.3}

\symup{H^{p}(X,\underline{\Omega}_X^{p})
=\operatorname{Ext}_{\Osh_X}^{p}(X;\Osh_X,\underline{\Omega}_X^{p})
\longleftarrow\operatorname{Ext}_{\Osh_X}^{p}(X;\Osh_Y,\underline{\Omega}_X^{p})
\xrightarrow{\ \sim\ }
H^{0}\!\left(X,\Extsh_{\Osh_X}^{p}(\Osh_Y,\underline{\Omega}_X^{p})\right)}
\end{equation}

Appelons \uline{cycle non singulier} de dimension $\symup{n-p}$ un élément du groupe
abélien libre engendré par les sous-variétés irréductibles non singulières de
dimension $\symup{n-p}$ dans $\symup{X}$. Alors la fonction $\symup{Y\longrightarrow P(Y)}$ se prolonge en
un homomorphisme du groupe des cycles non singuliers de dimension $\symup{n-p}$ sur $\symup{X}$,
dans le groupe $\symup{H^{p}(X,\underline{\Omega}_X^{p})}$.

Soient $\symup{Z^{n-p}}$ et $\symup{Z'^{n-p'}}$ des cycles non singuliers de dimension $\symup{n-p}$ et
$\symup{n-p'}$, on dit qu'ils se \uline{coupent transversalement}, si toute composante
de $\symup{Z}$ coupe transversalement toute composante de $\symup{Z'}$. Alors le cycle $\symup{Z.Z'}$ est
défini, c'est un cycle non singulier de dimension $\symup{n-p-p'}$. Ceci dit, on a le

\src{179}
\result{THÉORÈME}{1} \uline{Si} $\symup{Z^{n-p}}$ \uline{et} $\symup{Z'^{n-p'}}$ \uline{sont des cycles non singuliers
qui se coupent transversalement, alors on a}
\begin{equation}\tag{4.4}
  \symup{P_X(Z.Z')=P_X(Z).P_X(Z')}
\end{equation}
\uline{où le produit du deuxième membre est le cup-produit}
\[
  \symup{H^{p}(X,\underline{\Omega}_X^{p})\times H^{p'}(X,\underline{\Omega}_X^{p'})
  \longrightarrow H^{p+p'}(X,\underline{\Omega}_X^{p+p'}).}
\]
(On suppose que $\symup{X}$ est isomorphe à une partie localement fermée d'un espace
projectif)

Cette dernière hypothèse nous sert uniquement pour pouvoir conclure que tout
faisceau algébrique cohérent sur $\symup{X}$ est quotient d'un faisceau algébrique
cohérent localement libre (SERRE), donc admet une résolution gauche par des
faisceaux localement libres. Pour prouver le théorème 1, on peut supposer que
$\symup{Z}$ et $\symup{Z'}$ sont des sous-variétés irréductibles non singulières $\symup{Y}$ et $\symup{Y'}$ se
coupant transversalement. Soit $\symup{L_{*}}$ une résolution gauche de $\symup{\Osh_Y}$ par des
faisceaux localement libres, alors en vertu de la proposition 3, le diagramme
d'homomorphismes (4.3) s'identifie au diagramme
\[
\symup{R^{p}\Gamma(\underline{\Omega}_X^{p})
\xleftarrow{\alpha}
(\underline{\underline{R}}^{p}\Gamma)\!\left(\Homsh_{\Osh_X}(L_{*},\underline{\Omega}_X^{p})\right)
\xrightarrow{\overset{\beta}{\sim}}
\Gamma\!\left(H^{p}\bigl(\Homsh_{\Osh_X}(L_{*},\underline{\Omega}_X^{p})\bigr)\right)}
\]
où $\symup{\Gamma}$ est le foncteur \fquote{groupe des sections} sur la catégorie des
faisceaux abéliens sur $\symup{X}$, et $\symup{\underline{\underline{R}}^{p}\Gamma}$ désigne
l'hypercohomologie de dimension $\symup{p}$ dudit foncteur, et $\symup{R^{p}\Gamma}$ son
$\symup{p}$-ième foncteur dérivé. On supposera pour simplifier que $\symup{L_0=\Osh_X}$ et que
l'augmentation $\symup{L_0\longrightarrow\Osh_Y}$ est l'homomorphisme naturel (ce qui
est loisible), alors $\symup{\alpha}$ est déduit de l'homomorphisme de complexes
$\symup{\Osh_X\longrightarrow L}$ ($\symup{\Osh_X}$ étant considéré comme un complexe réduit
au degré $\symup{0}$), compte tenu que $\symup{\underline{\underline{R}}^{p}\Gamma(K)=R^{p}\Gamma(K_0)}$ si $\symup{K}$
est un complexe de faisceaux réduit au degré $\symup{0}$. L'homomorphisme $\symup{\beta}$ est un
\fquote{edge-homomorphism} bien connu. Considérons un diagramme analogue, relatif à
une résolution localement libre $\symup{L'_{*}}$ de $\symup{\Osh_{Y'}}$, et considérons le
diagramme commutatif d'accouplements :
\begin{equation}\tag{4.5}
\resizebox{\textwidth}{!}{$
\left\{\begin{array}{ccc}\symup{R^{p}\Gamma(\underline{\Omega}_X^{p})}&\symup{\longleftarrow
\underline{\underline{R}}^{p}\!\Gamma\left(\Homsh_{\Osh_X}(L_{*},\underline{\Omega}_X^{p})\right)}&\symup{\xrightarrow{\sim}
\Gamma\bigl(H^{p}\!\left(\Homsh_{\Osh_X}(L_{*},\underline{\Omega}_X^{p})\right)\bigr)}\\[3pt]\symup{\times}&\symup{\times}&\symup{\times}\\[3pt]\symup{R^{p'}\Gamma(\underline{\Omega}_X^{p'})}&\symup{\longleftarrow
\underline{\underline{R}}^{p'}\Gamma\!\left(\Homsh_{\Osh_X}(L'_{*},\underline{\Omega}_X^{p'})\right)}&\symup{\xrightarrow{\sim}
\Gamma\bigl(H^{p'}\!\left(\Homsh_{\Osh_X}(L'_{*},\underline{\Omega}_X^{p'})\right)\bigr)}\\[3pt]\symup{\downarrow}&\symup{\downarrow}&\symup{\downarrow}\\[3pt]\symup{R^{p+p'}\Gamma(\underline{\Omega}_X^{p+p'})}&\symup{\longleftarrow
\underline{\underline{R}}^{p+p'}\Gamma\!\left(
\Homsh_{\Osh_X}(L_{*}\otimes L'_{*},\underline{\Omega}_X^{p+p'})\right)}&\symup{\xrightarrow{\sim}
\Gamma\bigl(H^{p+p'}\!\left(
\Homsh_{\Osh_X}(L_{*}\otimes L'_{*},\underline{\Omega}_X^{p+p'})\right)\bigr)}\end{array}\right.
$}
\end{equation}

\src{180}
Les accouplements des deux colonnes de droite sont déduits de l'accouplement de
complexes de faisceaux
\[
\symup{\Homsh_{\Osh_X}(L_{*},\underline{\Omega}_X^{p})
\times\Homsh_{\Osh_X}(L'_{*},\underline{\Omega}_X^{p'})
\longrightarrow
\Homsh_{\Osh_X}(L_{*}\otimes L'_{*},\underline{\Omega}_X^{p+p'})}
\]
que l'on définit à l'aide du produit extérieur
$\symup{\underline{\Omega}_X^{p}\times\underline{\Omega}_X^{p'}\longrightarrow\underline{\Omega}_X^{p+p'}}$.
L'accouplement de la première colonne est le cup-produit (relatif au produit
extérieur). Je dis que la dernière ligne de (4.5) s'identifie au diagramme
d'homomorphismes analogue à (4.3), où $\symup{Y}$ est remplacé par $\symup{Y\cap Y'}$ et $\symup{p}$ par
$\symup{p+p'}$. Pour ceci, il suffit de montrer que $\symup{L\otimes L'}$ est une résolution
(évidemment localement libre) de $\symup{\Osh_{Y\cap Y'}}$. Or on a en effet
\[
  \symup{H_0(L\otimes L')=\Osh_Y\otimes\Osh_{Y'}=\Osh_{Y\cap Y'}}
\]
et
\[
  \symup{H_i(L\otimes L')=\operatorname{Tor}_{i}^{\Osh_X}(\Osh_Y,\Osh_{Y'})=0}
\]
pour $\symup{i>0}$, du fait que $\symup{Y}$ et $\symup{Y'}$ se coupent transversalement. Le théorème 1 résulte
maintenant de la formule :
\begin{equation}\tag{4.6}
  \symup{s_Y.s_{Y'}=s_{Y.Y'}}
\end{equation}
(où le produit du premier membre est celui de la dernière colonne de (4.5)).
Cette formule (4.6), de nature purement locale, se vérifie sans difficultés en
prenant pour $\symup{L_{*}}$ et $\symup{L'_{*}}$ les résolutions envisagées dans la proposition
4. On démontre de même (par une démonstration plus facile) que
$\symup{Z\longrightarrow P_X(Z)}$ est compatible avec le produit cartésien :
\begin{equation}\tag{4.7}
  \symup{P_{X\times X'}(Z\times Z')=P_X(Z)\otimes P_{X'}(Z')}
\end{equation}
(formule valable si $\symup{Z}$ resp. $\symup{Z'}$, est un cycle non singulier sur la variété
non singulière $\symup{X}$, resp. $\symup{X'}$, $\symup{Z\times Z'}$ étant considéré comme un cycle non
singulier sur $\symup{X\times X'}$). De (4.4) et (4.7) on déduit que $\symup{P_X(Z)}$ est aussi
compatible avec l'opération \fquote{image inverse} par un morphisme de variétés non
singulières $\symup{f:X\longrightarrow X'}$ :
\begin{equation}\tag{4.8}
  \symup{P_X\bigl(f^{-1}(Z')\bigr)=f^{*}\bigl(P_{X'}(Z')\bigr)}
\end{equation}
formule valable si $\symup{Z'}$ est un cycle non singulier sur $\symup{X'}$ tel que $\symup{f}$ soit
\fquote{transversale} à $\symup{Z'}$, i.e. tel que le graphe de $\symup{f}$ soit transversal au cycle
$\symup{X\times Z'}$ dans $\symup{X\times X'}$.

\src{181}
\result{COROLLAIRE}{1} \uline{Soient} $\symup{X}$, $\symup{X'}$ \uline{deux variétés non singulières
localement fermées dans un espace projectif, supposons} $\symup{X'}$ \uline{complète. Soient} $\symup{U}$
\uline{un cycle non singulier sur} $\symup{X\times X'}$, $\symup{a}$ \uline{et} $\symup{b}$ \uline{deux points de} $\symup{X'}$ \uline{tels que}
$\symup{U}$ \uline{coupe transversalement les cycles} $\symup{X\times(a)}$ \uline{et} $\symup{X\times(b)}$, \uline{soient} $\symup{Z}$
\uline{et} $\symup{Z'}$ \uline{les cycles non singuliers sur} $\symup{X}$ \uline{tels que}
$\symup{Z\times(a)=(X\times(a)).U}$, $\symup{Z'\times(b)=(X\times(b)).U}$.
\uline{Alors on a}
\[
  \symup{P_X(Z)=P_X(Z').}
\]

En effet, soit $\symup{f_a:X\longrightarrow X\times X'}$ défini par
$\symup{f_a(x)=(x,a)}$, on a alors en vertu de (4.8) la formule
$\symup{P(Z)=f_a^{*}(P(U))}$, de même $\symup{P(Z')=f_b^{*}(P(U))}$. Or, utilisant la formule de
Künneth
\[
  \symup{H^{*}(X\times X',\underline{\Omega}_{X\times X'}^{*})
  =H^{*}(X,\underline{\Omega}_X^{*})\otimes H^{*}(X',\underline{\Omega}_{X'}^{*})}
\]
et le fait que $\symup{H^{0}(X',\Osh_{X'})}$ est réduit aux scalaires, on obtient
facilement $\symup{f_a^{*}=f_b^{*}}$, d'où la conclusion.

Pour tout $\symup{x\in X}$, $\symup{(x)}$ est une sous-variété non singulière de $\symup{X}$ de
codimension $\symup{n}$, et définit donc un élément $\symup{\varepsilon_x}$ de
$\symup{H^{n}(X,\underline{\Omega}_X^{n})}$. Si $\symup{X}$ est une variété projective non singulière, il
résulte du corollaire 1 que $\symup{\varepsilon_x}$ ne dépend pas du point $\symup{x}$ choisi, on
le note $\symup{\varepsilon_X}$ et on l'appelle \uline{classe fondamentale} de
$\symup{H^{n}(X,\underline{\Omega}_X^{n})}$.

\result{REMARQUE}{} Pour avoir une théorie satisfaisante, il faudrait définir
$\symup{P_X(Z)}$ pour des cycles $\symup{Z}$ quelconques, et prouver le théorème 1 pour une
intersection propre de cycles. (Au moment d'écrire cet exposé, cela n'est pas
encore fait en toute généralité). Admettant que cela est fait, le corollaire 1
devient : si $\symup{Z}$ et $\symup{Z'}$ sont deux cycles algébriquement équivalents, alors
$\symup{P_X(Z)=P_X(Z')}$ (énoncé qui ne semble pas résulter de ce qui précède, même si
$\symup{Z}$ et $\symup{Z'}$ sont non singuliers).

\fsection{\normalfont5}{\normalfont\uline{Le théorème de dualité}.}

Dans ce numéro, $\symup{X}$ désigne une variété projective non singulière de dimension
$\symup{n}$.

\result{THÉORÈME}{2} \uline{La classe fondamentale} $\symup{\varepsilon_X}$ \uline{de}
$\symup{H^{n}(X,\underline{\Omega}_X^{n})}$ \uline{est une base de cet espace vectoriel}.

(La démonstration sera donnée plus bas). Moyennant le théorème précédent on peut
donc identifier $\symup{H^{n}(X,\underline{\Omega}_X^{n})}$ au corps $\symup{\underline{k}}$. Considérons maintenant les
accouplements envisagés au n\textsuperscript{o} 2, qui donnent en particulier un
accouplement

\src{182}
\[
\symup{\operatorname{Ext}_{\Osh_X}^{p}(X;\Osh_X,F)
\times\operatorname{Ext}_{\Osh_X}^{n-p}(X;F,\underline{\Omega}_X^{n})
\longrightarrow
\operatorname{Ext}_{\Osh_X}^{n}(X;\Osh_X,\underline{\Omega}_X^{n}),\quad\text{i.e.}}
\]
\begin{equation}\tag{5.1}
  \symup{H^{p}(X,F)
  \times\operatorname{Ext}_{\Osh_X}^{n-p}(X;F,\underline{\Omega}_X^{n})
  \longrightarrow H^{n}(X,\underline{\Omega}_X^{n})}
\end{equation}
Compte tenu du théorème 2, cet accouplement définit un homomorphisme
\begin{equation}\tag{5.2}
  \symup{\operatorname{Ext}_{\Osh_X}^{n-p}(X;F,\underline{\Omega}_X^{n})
  \longrightarrow\bigl(H^{p}(X,F)\bigr)'}
\end{equation}
Cet homomorphisme est fonctoriel en $\symup{F}$, et permute aux homomorphismes cobords
relatifs aux suites exactes $\symup{0\longrightarrow F'\longrightarrow F
\longrightarrow F''\longrightarrow0}$.

\result{THÉORÈME}{3} \uline{L'homomorphisme} (5.2) \uline{est un isomorphisme}.

On retrouve en particulier le résultat de Serre :

\result{COROLLAIRE}{} \uline{Soient} $\symup{E}$ \uline{un fibré vectoriel algébrique sur} $\symup{X}$,
$\symup{\Osh_X(E)}$ \uline{le faisceau des germes de sections régulières de} $\symup{E}$, \uline{on a alors des
isomorphismes canoniques} :
\begin{equation}\tag{5.3}
  \symup{\bigl(H^{p}(X,\Osh_X(E))\bigr)'
  =H^{n-p}\bigl(X,\underline{\Omega}_X^{n}\otimes\Osh_X(E')\bigr)}
\end{equation}
Il suffit d'appliquer le théorème 3 et le corollaire 1 à la proposition 1.

Les théorèmes 2 et 3 vont résulter de l'énoncé suivant :

\noindent(D) \uline{L'homomorphisme suivant}
\begin{equation}\tag{5.2 bis}
  \symup{\operatorname{Ext}_{\Osh_X}^{n-p}(X;F,\underline{\Omega}_X^{n})
  \longrightarrow\bigl(H^{p}(X,F)\bigr)'\otimes L
  \qquad (L=H^{n}(X,\underline{\Omega}_X^{n})),}
\end{equation}
\uline{déduit de l'accouplement (5.1) est un isomorphisme}.

Montrons en effet que (D) implique le théorème 2. Soit
$\symup{\underline{k}_x=\Osh_{(x)}}$ le faisceau des anneaux locaux de la variété réduite au point
$\symup{x\in X}$, considérons l'homomorphisme canonique
$\symup{\Osh_X\longrightarrow \underline{k}_x}$ et l'homomorphisme associé
\begin{equation}\tag{5.3 bis}
  \symup{H^{0}(X,\Osh_X)\longrightarrow H^{0}(X,\underline{k}_x)}
\end{equation}
son transposé s'identifie à l'homomorphisme
\begin{equation}\tag{5.4}
  \symup{\operatorname{Ext}_{\Osh_X}^{n}(X;\underline{k}_x,\underline{\Omega}_X^{n})\otimes L'
  \longrightarrow
  \operatorname{Ext}_{\Osh_X}^{n}(X;\Osh_X,\underline{\Omega}_X^{n})\otimes L'}
\end{equation}
déduit de l'homomorphisme sur les $\symup{\operatorname{Ext}^{n}}$ associé à
$\symup{\Osh_X\longrightarrow \underline{k}_x}$ :
\begin{equation}\tag{5.5}
  \symup{\operatorname{Ext}_{\Osh_X}^{n}(X;\underline{k}_x,\underline{\Omega}_X^{n})
  \longrightarrow
  \operatorname{Ext}_{\Osh_X}^{n}(X;\Osh_X,\underline{\Omega}_X^{n})
  =H^{n}(X,\underline{\Omega}_X^{n})}
\end{equation}

\src{183}
Comme (5.3 bis) est un isomorphisme,\footnote{Le texte imprimé donne une
seconde formule « (5.3) » à la page 182 et y renvoie ici ; l'identifiant
« (5.3 bis) » distingue cette formule de la première (5.3).} il en est de même
de (5.4), donc aussi de (5.5). Comme $\symup{s_{(x)}}$ est une base de
$\symup{\operatorname{Ext}_{\Osh_X}^{n}(X;\underline{k}_x,\underline{\Omega}_X^{n})}$ en vertu de (4.2), il
s'ensuit bien que son image $\symup{\varepsilon_X}$ est une base de
$\symup{H^{n}(X,\underline{\Omega}_X^{n})}$.

Reste à prouver l'énoncé (D), qui résultera de façon purement formelle des faits
élémentaires résumés dans les lemmes suivants. On y suppose que $\symup{X}$ est une
partie fermée (singulière ou non) de l'espace projectif $\symup{P}$ de dimension $\symup{r}$.
On utilise la notation $\symup{\Osh_P(m)}$ pour le faisceau sur $\symup{P}$ noté $\symup{\Osh(m)}$ dans
[3], et la notation $\symup{\Osh_X(m)}$ pour le faisceau analogue sur $\symup{X}$.

\result{LEMME}{2} \uline{L'énoncé} (D) \uline{est vrai si} $\symup{X=P}$ \uline{et} $\symup{F=\Osh_P(m)}$

Ce lemme se vérifie par un calcul direct. Le calcul explicite des
$\symup{H^{i}(P,\Osh_P(m))}$ se trouve dans [3], mais il peut se faire plus
élémentairement. La computation du cup-produit $\symup{H^{i}(P,\Osh_P(m))\times H^{j}(P,\Osh_P(m'))\longrightarrow H^{i+j}(P,\Osh_P(m+m'))}$ nécessaire pour expliciter l'accouplement (5.1) n'offre pas de
difficultés.\footnote{Le texte imprimé porte $\symup{\Osh_P(r+r')}$ ; les deux facteurs
$\symup{\Osh_P(m)}$ et $\symup{\Osh_P(m')}$ imposent $\symup{\Osh_P(m+m')}$.}

\result{LEMME}{3} \uline{Tout faisceau algébrique cohérent} $\symup{F}$ \uline{sur} $\symup{X}$ \uline{est isomorphe à
un faisceau quotient d'un faisceau} $\symup{\Osh_X(-m)^{k}}$, \uline{où on peut supposer} $\symup{m}$
\uline{aussi grand que l'on veut}.

Résulte du fait que $\symup{F\otimes\Osh_X(m)}$ est \fquote{engendré par ses sections} pour
$\symup{m}$ grand, cf. [3].

\result{LEMME}{4} \uline{Soit} $\symup{i>0}$. \uline{Alors}
$\symup{H^{r-i}(P,\Osh_P(-m))=0}$ \uline{pour} $\symup{m}$ \uline{grand ; et pour tout faisceau algébrique
cohérent} $\symup{B}$ \uline{sur} $\symup{X}$, \uline{on a} $\symup{\operatorname{Ext}_{\Osh_X}^{i}(X;\Osh_X(-m),B)=0}$ \uline{pour} $\symup{m}$ \uline{grand}.

Le premier énoncé résulte des calculs explicites mentionnés plus haut, pour le
deuxième on note qu'on a un isomorphisme
\[
  \symup{\operatorname{Ext}_{\Osh_X}^{i}(X;\Osh_X(-m),B)
  =H^{i}(X,B\otimes\Osh(m))}
\]
(proposition 1, corollaire 1), d'où la conclusion en vertu d'un résultat bien
connu de [3]. Conjuguant les deux lemmes précédents, on trouve le

\result{COROLLAIRE}{} \uline{Soit} $\symup{i>0}$. \uline{Alors le foncteur}
$\symup{F\longrightarrow H^{r-i}(P,F)}$ \uline{sur la catégorie des faisceaux algébriques cohérents
sur} $\symup{P}$ \uline{est coeffaçable, et il en est de même du foncteur}
$\symup{\operatorname{Ext}_{\Osh_X}^{i}(X;F,B)}$ \uline{sur la catégorie des faisceaux
algébriques cohérents sur} $\symup{X}$.

\src{184}
\result{LEMME}{5} \uline{Soient} $\symup{A}$ \uline{et} $\symup{B}$ \uline{deux faisceaux algébriques cohérents sur}
$\symup{X}$, \uline{posons} $\symup{A(m)=A\otimes\Osh_X(m)}$. \uline{Alors pour} $\symup{m}$ \uline{assez grand,
l'homomorphisme canonique}
\[
  \symup{\operatorname{Ext}_{\Osh_X}^{i}(X;A(-m),B)
  \longrightarrow H^{0}\bigl(X,\Extsh_{\Osh_X}^{i}(A(-m),B)\bigr)
  =H^{0}\bigl(X,\Extsh_{\Osh_X}^{i}(A,B)(m)\bigr)}
\]
\uline{est un isomorphisme}.

Cela résulte aussitôt de la suite spectrale de la proposition 1 appliquée à
$\symup{A(-m)}$ et $\symup{B}$, puisque l'on aura
\[
  \symup{E_{2}^{p,q}(A(-m),B)
  =H^{p}\bigl(X,\Extsh_{\Osh_X}^{q}(A(-m),B)\bigr)
  =H^{p}\bigl(X,\Extsh_{\Osh_X}^{q}(A,B)(m)\bigr),}
\]
qui est nul pour $\symup{p>0}$ et $\symup{m}$ grand.

Démontrons alors (D) dans le cas où $\symup{X=P}$. Nous prouvons d'abord que (5.2 bis)
est un isomorphisme pour $\symup{p=n}$ ; comme les deux membres sont alors des
foncteurs exacts à gauche (puisque $\symup{H^{r+1}(P,F)=0}$), il résulte du lemme 3
qu'il suffit de prouver l'assertion pour $\symup{F=\Osh_P(-m)}$, or elle est alors
contenue dans le lemme 2. Comme les homomorphismes (5.2 bis) sont fonctoriels et
compatibles avec les opérateurs cobords, et que pour $\symup{p<n}$ les deux membres de
(5.2 bis) sont des foncteurs coeffaçables en $\symup{F}$ (corollaire du lemme 4), il
s'ensuit alors par un raisonnement standard que (5.2 bis) est un isomorphisme
pour tout $\symup{p}$. Cela prouve le théorème de dualité pour l'espace projectif.

Supposons maintenant $\symup{X}$ quelconque, mais non singulier. D'après le théorème de
dualité pour $\symup{P}$, on a un isomorphisme,
\[
  \symup{H^{n}(X,F)'=H^{n}(P,F)'
  =\operatorname{Ext}_{\Osh_P}^{r-n}(P;F,\underline{\Omega}_P^{r})}
\]
En vertu du lemme 1 (n\textsuperscript{o} 4) le dernier membre s'identifie à
\[
  \symup{\operatorname{Hom}_{\Osh_P}(P;F,\underline{\omega}_X^{n})
  =\operatorname{Hom}_{\Osh_X}(X;F,\underline{\Omega}_X^{n})
  =\operatorname{Ext}_{\Osh_X}^{0}(X;F,\underline{\Omega}_X^{n})}
\]
d'où un isomorphisme
\begin{equation}\tag{5.6}
  \symup{H^{n}(X,F)'=\operatorname{Hom}_{\Osh_X}(X;F,\underline{\Omega}_X^{n})
  =\operatorname{Ext}_{\Osh_X}^{0}(X;F,\underline{\Omega}_X^{n})}
\end{equation}
Faisant $\symup{F=\underline{\Omega}_X^{n}}$, on trouve un isomorphisme
\begin{equation}\tag{5.7}
  \symup{\gamma:H^{n}(X,\underline{\Omega}_X^{n})\xrightarrow{\sim}\underline{k}}
\end{equation}

\src{185}
On vérifie que l'isomorphisme (5.6) n'est autre que (5.2 bis) pour $\symup{p=n}$, quand
on y fait $\symup{L=\underline{k}}$ grâce à (5.7). Par suite, (5.2 bis) est un isomorphisme pour
$\symup{p=n}$. Pour prouver que c'est un isomorphisme pour tout $\symup{p}$, il suffit encore
de prouver que pour $\symup{p<n}$, les deux membres de (5.2 bis) sont des foncteurs
coeffaçables en $\symup{F}$, et a fortiori (compte tenu du lemme 3) que les deux membres
sont nuls quand on fait $\symup{F=\Osh_X(-m)}$ avec $\symup{m}$ grand. Or pour le premier
membre cela est vrai en vertu du lemme 4, et pour le deuxième membre on écrit,
utilisant le théorème de dualité sur $\symup{P}$ :
\[
  \symup{H^{p}(X,\Osh_X(-m))'
  =\operatorname{Ext}_{\Osh_P}^{r-p}
    (P;\Osh_X(-m),\underline{\Omega}_P^{r}).}
\]
Le deuxième membre est nul pour $\symup{p<n}$ et $\symup{m}$ grand, comme il résulte du lemme 5
(où on fait $\symup{X=P}$) et du fait que $\symup{\Osh_X}$ est en tant que faisceau algébrique
cohérent sur $\symup{P}$, de dimension cohomologique $\symup{\leq r-n}$, (car $\symup{X}$ est non
singulière) d'où
\[
  \symup{\Extsh_{\Osh_P}^{r-p}(\Osh_X,\underline{\Omega}_P^{r})=0
  \qquad\text{pour }p<n.}
\]

\fsection{\normalfont6}{\normalfont\uline{Le théorème de dualité pour les variétés singulières}.}

Soit $\symup{X}$ une partie fermée de dimension $\symup{n}$ de l'espace projectif $\symup{P}$ de
dimension $\symup{r}$. La formule (5.6) doit s'écrire ici
\begin{equation}\tag{6.1}
  \symup{H^{n}(X,F)'
  \simeq\operatorname{Hom}_{\Osh_X}(X;F,\underline{\omega}_X^{n})
  =\operatorname{Ext}_{\Osh_X}^{0}(X;F,\underline{\omega}_X^{n})}
\end{equation}
où on a posé
\begin{equation}\tag{6.2}
  \symup{\underline{\omega}_X^{n}=E^{0}(\Osh_X)
  =\Extsh_{\Osh_P}^{r-n}(\Osh_X,\underline{\Omega}_P^{r})}
\end{equation}
Comme il a été dit dans la proposition 5, le faisceau ainsi introduit ne dépend
pas en fait de l'immersion choisie de $\symup{X}$ dans une variété non singulière $\symup{P}$.
Faisant, dans (6.1), $\symup{F=\underline{\omega}_X^{n}}$, on trouve
\begin{equation}\tag{6.2 bis}
  \symup{H^{n}(X,\underline{\omega}_X^{n})'
  \simeq\operatorname{Hom}_{\Osh_X}(X;\underline{\omega}_X^{n},\underline{\omega}_X^{n})}
\end{equation}
\footnote{Le texte imprimé numérote aussi cette formule « (6.2) » ; le suffixe
« bis » distingue les deux occurrences.} d'où l'existence d'un élément
privilégié dans $\symup{H^{n}(X,\underline{\omega}_X^{n})'}$, correspondant à l'homomorphisme
identique de $\symup{\underline{\omega}_X^{n}}$ dans lui-même :
\begin{equation}\tag{6.3}
  \symup{\gamma:H^{n}(X,\underline{\omega}_X^{n})\longrightarrow\underline{k}}
\end{equation}

\src{186}
Considérons alors les accouplements définis par la composition des $\symup{\operatorname{Ext}}$ :
\begin{equation}\tag{6.4}
  \symup{H^{p}(X,F)\times
  \operatorname{Ext}_{\Osh_X}^{n-p}(X;F,\underline{\omega}_X^{n})
  \longrightarrow H^{n}(X,\underline{\omega}_X^{n})}
\end{equation}
et composons-les avec l'homomorphisme $\symup{\gamma}$ de (6.3), on obtient des
homomorphismes fonctoriels, compatibles avec les opérateurs bord :
\begin{equation}\tag{6.5}
  \symup{\operatorname{Ext}_{\Osh_X}^{n-p}(X;F,\underline{\omega}_X^{n})
  \longrightarrow H^{p}(X,F)'}
\end{equation}
(généralisant (5.2)). On vérifie que pour $\symup{p=n}$, on obtient ainsi
l'isomorphisme (6.1). Ceci posé, on a le

\result{THÉORÈME}{3 bis} \uline{Les quatre conditions suivantes sur} $\symup{X}$ \uline{sont
équivalentes pour un entier} $\symup{k\geq0}$ \uline{donné} :
\begin{enumerate}
  \item[\textnormal{i.}] \uline{L'homomorphisme fonctoriel} (6.5) \uline{est un isomorphisme pour}
    $\symup{n-k\leq p\leq n}$.
  \item[\textnormal{ii.}] \uline{On a} $\symup{H^{p}(X,\Osh_X(-m))=0}$ \uline{pour} $\symup{m}$ \uline{grand et pour}
    $\symup{n-k\leq p<n}$.
  \item[\textnormal{iii.}] \uline{Le foncteur} $\symup{H^{p}(X,F)}$ \uline{sur la catégorie des faisceaux
    algébriques cohérents sur} $\symup{X}$ \uline{est coeffaçable pour} $\symup{n-k\leq p<n}$.
  \item[\textnormal{iv.}] \uline{On a}
    $\symup{E^{i}(\Osh_X)=\Extsh_{\Osh_P}^{r-n+i}(\Osh_X,\underline{\omega}_P^{r})=0}$
    \uline{pour} $\symup{0<i\leq k}$.
\end{enumerate}

\noindent DÉMONSTRATION. -- \textnormal{i.}$\symup{\Rightarrow}$\textnormal{ii.} en
vertu du lemme 4, \textnormal{ii.}$\symup{\Rightarrow}$\textnormal{iii.} en vertu du lemme 3,
\textnormal{iii.}$\symup{\Rightarrow}$\textnormal{i.} par un raisonnement standard bien connu,
compte tenu que les deux membres de (6.5) sont alors des foncteurs coeffaçables
pour $\symup{n-k\leq p<n}$ (le premier l'étant en vertu du lemme 4). Enfin prouvons
\textnormal{ii.}$\symup{\Leftrightarrow}$\textnormal{iv.} Cela résulte du corollaire à la
proposition qui suit :

\result{PROPOSITION}{6} \uline{Soit} $\symup{F}$ \uline{un faisceau algébrique cohérent sur} $\symup{X}$, \uline{et
soit} $\symup{i}$ \uline{un entier. Alors pour} $\symup{m}$ \uline{assez grand, on a un isomorphisme} :
\begin{equation}\tag{6.6}
  \symup{H^{i}(X,F(-m))'
  \simeq H^{0}\bigl(X,E^{n-i}(F)(m)\bigr),}
\end{equation}
\uline{où on pose} (comparer n\textsuperscript{o} 3, proposition 5) :
\begin{equation}\tag{6.7}
  \symup{E^{j}(F)=\Extsh_{\Osh_P}^{r-n+j}(F,\underline{\Omega}_P^{r})}
\end{equation}
En effet, en vertu du théorème de dualité sur $\symup{P}$, le premier membre de (6.6)
est isomorphe à
$\symup{\operatorname{Ext}_{\Osh_P}^{r-i}(P;F(-m),\underline{\Omega}_P^{r})}$, donc (6.6) résulte
du lemme 5.

\src{187}
\result{COROLLAIRE}{} \uline{Pour qu'on ait} $\symup{H^{i}(X,F(-m))=0}$ \uline{pour} $\symup{m}$ \uline{grand, il
faut et il suffit que} $\symup{E^{n-i}(F)=0}$.

Rappelons que les $\symup{E^{j}(F)}$ ne dépendent pas en fait de l'immersion projective
envisagée. La condition du corollaire est de nature purement locale, par suite,
si elle est vérifiée pour $\symup{F}$, elle l'est pour tout faisceau localement
isomorphe à un faisceau $\symup{F^{n}}$. En particulier, si cette condition est vérifiée
pour $\symup{\Osh_X}$, elle l'est pour tout faisceau algébrique cohérent localement
libre. C'est par exemple le cas pour tout $\symup{i<n}$ si $\symup{X}$ est non singulière, pour
$\symup{i=0}$ si aucune composante de $\symup{X}$ n'est réduite à un seul point, pour $\symup{i=0,1}$
si $\symup{X}$ est normale et toutes ses composantes sont de dimension $\symup{>1}$ (cf.
[3]).\footnote{Le texte imprimé porte ici \fquote{$\symup{S}$ est normale} ; le paragraphe
ne fait intervenir que la variété $\symup{X}$.} Pour qu'il en soit ainsi pour tout
$\symup{i<k}$, il faut et il suffit par définition que les anneaux locaux
$\symup{\Osh_x}$ ($\symup{x\in X}$) soient de \fquote{codimension homologique} $\symup{\geq k}$ (cf [4]
pour cette notion). Si $\symup{k=n}$, cela signifie en vertu du théorème 3 bis que le
théorème de dualité est vrai pour $\symup{X}$, i.e. que (6.5) est un isomorphisme pour
tout $\symup{p}$ et tout $\symup{F}$. On peut donner de nombreuses conditions équivalentes sur
les anneaux locaux $\symup{\Osh_x}$ pour qu'il en soit ainsi (NAGATA), par exemple celle
de vérifier le théorème d'équidimensionnalité de Cohen-Macaulay. Il en est ainsi
par exemple si localement $\symup{X}$ est une \fquote{intersection complète} dans une
variété ambiante non singulière.

\fsection{\normalfont7}{\normalfont\uline{La dualité de Poincaré}.}

Soit $\symup{X}$ une variété projective non singulière de dimension $\symup{n}$. Alors
$\symup{H^{**}(X)=H^{*}(X,\underline{\Omega}_X^{*})}$ est une algèbre anticommutative bigraduée de dimension finie, que nous graduons
par le degré total, $\symup{H^{p,q}(X)=H^{p}(X,\underline{\Omega}_X^{q})}$ étant donc de degré
$\symup{p+q}$. Les degrés de $\symup{H^{*}(X)}$ sont compris entre $\symup{0}$ et $\symup{2n}$. En vertu du
théorème 2 et du corollaire au théorème 3, $\symup{H^{*}(X)}$ est une \fquote{algèbre de
Poincaré} de dimension $\symup{2n}$, i.e. $\symup{H^{2n}(X)}$ est muni d'un isomorphisme avec
le corps de base $\symup{\underline{k}}$, et le produit
$\symup{H^{m}(X)\times H^{2n-m}(X)\longrightarrow H^{2n}(X)=\underline{k}}$ est une dualité entre $\symup{H^{m}(X)}$ et $\symup{H^{2n-m}(X)}$. De plus, si $\symup{Y}$ est une
autre variété projective non singulière, la formule de Künneth pour les
faisceaux algébriques cohérents donne
\begin{equation}\tag{7.1}
  \symup{H^{*}(X\times Y)=H^{*}(X)\otimes H^{*}(Y)}
\end{equation}
isomorphisme compatible avec les structures d'algèbres de Poincaré. De plus,
$\symup{H^{*}(X)}$ est, en tant qu'algèbre commutative, un foncteur contravariant en
$\symup{X}$, un morphisme $\symup{f:Y\longrightarrow X}$ définissant de façon évidente un
homomorphisme d'algèbres graduées

\src{188}
\begin{equation}\tag{7.2}
  \symup{f^{*}:H^{*}(X)\longrightarrow H^{*}(Y)}
\end{equation}
Comme il s'agit d'algèbres de Poincaré, on obtient par transposition un
homomorphisme d'espaces vectoriels :
\begin{equation}\tag{7.3}
  \symup{f_{*}:H^{*}(Y)\longrightarrow H^{*}(X)}
\end{equation}
On a vu au n\textsuperscript{o} 4 que l'effet de $\symup{f^{*}}$ sur les classes de
cohomologie correspondant à des cycles non singuliers s'interprète
géométriquement en prenant les classes de cohomologie qui correspondent à leurs
images inverses. Il importe dans le cas actuel de montrer que (7.3) correspond
de même à l'opération \fquote{image directe} de cycles. Cela résulte, sous des
conditions de non singularité convenables du moins, du cas particulier suivant :

\result{THÉORÈME}{4} \uline{Si} $\symup{f}$ \uline{est l'application identique d'une sous-variété non
singulière} $\symup{Y^{m}}$ \uline{de} $\symup{X^{n}}$ \uline{dans} $\symup{X^{n}}$, \uline{alors, désignant par} $\symup{1_Y}$
\uline{l'élément unité de} $\symup{H(Y)}$ \uline{on a}
\begin{equation}\tag{7.4}
  \symup{f_{*}(1_Y)=P_X(Y)}
\end{equation}
\uline{où le deuxième membre est la classe de cohomologie dans} $\symup{X}$ \uline{associée à} $\symup{Y}$.

Cette formule est équivalente à
\begin{equation}\tag{7.4 bis}
  \symup{\bigl\langle\xi^{m,m},P_X(Y)\bigr\rangle\varepsilon_Y
  =f^{*}(\xi^{m,m})
  \qquad
  \bigl(\xi^{m,m}\in H^{m}(X,\underline{\Omega}_X^{m})\bigr)}
\end{equation}
\footnote{Le texte imprimé porte $\symup{f_{*}(\xi^{m,m})}$. Puisque
$\symup{\xi^{m,m}\in H^{m}(X,\underline{\Omega}_X^{m})}$ et que les deux membres appartiennent à
$\symup{H^{m}(Y,\underline{\Omega}_Y^{m})}$, l'application requise est $\symup{f^{*}}$.} où
$\symup{\varepsilon_Y}$ est l'élément fondamental dans $\symup{H^{m}(Y,\underline{\Omega}_Y^{m})}$, et
constitue, dans le cas des variétés projectives non singulières, une nouvelle
définition de la classe de cohomologie associée à $\symup{Y}$. Pour démontrer le
théorème 4, on interprète grâce au théorème 3 la transposée de l'homomorphisme
\[
  \symup{H^{m}(X,\underline{\Omega}_X^{m})
  \longrightarrow H^{m}(Y,\underline{\Omega}_Y^{m})
  =H^{m}(X,\underline{\Omega}_Y^{m})}
\]
comme l'homomorphisme
\begin{equation}\tag{7.5}
\symup{H^{n-m}(X,\underline{\Omega}_X^{n-m})
  \xrightarrow{\sim}
  \operatorname{Ext}_{\Osh_X}^{n-m}(X;\underline{\Omega}_X^{m},\underline{\Omega}_X^{n})
  \longleftarrow
  \operatorname{Ext}_{\Osh_X}^{n-m}(X;\underline{\Omega}_Y^{m},\underline{\Omega}_X^{n})
  \xrightarrow{\sim}
  \operatorname{Hom}_{\Osh_X}(X;\underline{\Omega}_Y^{m},\underline{\Omega}_Y^{m})}
\end{equation}
On vérifie que l'élément $\symup{1_Y}$ du dual de $\symup{H^{m}(Y,\underline{\Omega}_Y^{m})}$ s'identifie
à l'élément du membre de droite correspondant à l'endomorphisme identique de
$\symup{\underline{\Omega}_Y^{m}}$, et d'autre part que l'image de cet élément dans
$\symup{H^{n-m}(X,\underline{\Omega}_X^{n-m})}$ est bien $\symup{P_X(Y)}$. Ces compatibilités, qui auraient
pu être données dès le n\textsuperscript{o} 4, peuvent s'énoncer, et sont
valables en fait, pour des variétés non singulières quelconques, la deuxième

\src{189}
par exemple résultant de la commutativité du diagramme suivant d'homomorphismes
canoniques :
\begin{equation}\tag{7.6}
\resizebox{\textwidth}{!}{$
\left\{
\begin{array}{ccccc}
\symup{\operatorname{Ext}_{\Osh_X}^{n-m}(X;\underline{\Omega}_X^{m},\underline{\Omega}_X^{n})}
&
\symup{\longleftarrow}
&
\symup{\operatorname{Ext}_{\Osh_X}^{n-m}(X;\underline{\Omega}_Y^{m},\underline{\Omega}_X^{n})}
&
\symup{\xrightarrow{\sim}}
&
\symup{\operatorname{Hom}_{\Osh_X}(X;\underline{\Omega}_Y^{m},\underline{\Omega}_Y^{m})}
\\[3pt]
\symup{\Big\downarrow}
&&&&
\symup{\Big\downarrow}
\\[3pt]
\symup{\operatorname{Ext}_{\Osh_X}^{n-m}(X;\Osh_X,\underline{\Omega}_X^{n-m})}
&
\symup{\longleftarrow}
&
\symup{\operatorname{Ext}_{\Osh_X}^{n-m}(X;\Osh_Y,\underline{\Omega}_X^{n-m})}
&
\symup{\xrightarrow{\sim}}
&
\symup{\operatorname{Hom}_{\Osh_X}(X;\underline{\Omega}_X^{m},\underline{\Omega}_Y^{m})}
\end{array}
\right.
$}
\end{equation}

On obtient ainsi un équivalent exact du formalisme de la dualité de Poincaré
sur les variétés orientées compactes. En particulier, le théorème 4 permet de
déterminer la classe de cohomologie associée à la diagonale de $\symup{X\times X}$. Par
un raisonnement bien connu, on en déduit par exemple une \uline{formule de Lefschetz} :

\result{THÉORÈME}{5} \uline{Soit} $\symup{f}$ \uline{un endomorphisme d'une variété projective non
singulière} $\symup{X}$, \uline{telle que les points fixes de} $\symup{f}$ \uline{soient de multiplicité} $\symup{1}$.
\uline{Alors le nombre de ces points fixes est congru}, $\symup{\operatorname{mod}}$ \uline{la caractéristique de}
$\symup{\underline{k}}$ \uline{à la somme alternée des traces des endomorphismes des} $\symup{H^{i}(X)}$ \uline{définis
par} $\symup{f}$.

La restriction sur $\symup{f}$ qu'on a dû faire tient aux difficultés mentionnées dans
la remarque du n\textsuperscript{o} 4. On notera, toutefois, que la formule de
Lefschetz reste vraie si $\symup{f}$ est \fquote{homotope} à un endomorphisme dont tous les
points fixes sont de multiplicité $\symup{1}$.

\fsection{\normalfont8}{\normalfont\uline{Généralisation du théorème de dualité}.}

Soit $\symup{X}$ une variété algébrique non singulière telle que tout faisceau
algébrique cohérent $\symup{F}$ sur $\symup{X}$ soit isomorphe à un quotient d'un faisceau
algébrique cohérent localement libre (ce qui est le cas si $\symup{X}$ est localement
fermé dans un espace projectif). Alors tout faisceau algébrique cohérent $\symup{F}$ sur
$\symup{X}$ admet une résolution finie $\symup{L}$ par des faisceaux localement libres, et pour
deux telles résolutions on peut toujours trouver une troisième, et un
homomorphisme de celle-ci dans les deux premières compatibles avec les
augmentations. De même, si $\symup{L}$ est une résolution finie localement libre de
$\symup{F}$, et si on se donne un homomorphisme $\symup{F'\longrightarrow F}$, alors il existe
une résolution finie localement libre $\symup{L'}$ de $\symup{F'}$ et un homomorphisme
$\symup{L'\longrightarrow L}$ compatible avec $\symup{F'\longrightarrow F}$, qu'on peut
supposer surjectif si $\symup{F'\longrightarrow F}$ l'est. Cela permet de définir, étant
donnés deux entiers $\symup{r,s\geq0}$, deux multifoncteurs, cohomologiques, en des
arguments $\symup{A_1,\ldots,A_r}$ ; $\symup{B_1,\ldots,B_s}$

\src{190}
variant dans la catégorie des faisceaux algébriques cohérents sur $\symup{X}$, et dont
les valeurs sont dans la catégorie des faisceaux algébriques cohérents sur $\symup{X}$,
resp. dans la catégorie des modules sur $\symup{H^{0}(X,\Osh_X)}$, par les formules :
\begin{equation}\tag{8.1}
\resizebox{\textwidth}{!}{$
\left\{
\begin{aligned}
\symup{\underline{T}_{r}^{s*}(A_1,\ldots,A_r;B_1,\ldots,B_s)}
&
\symup{={}H^{*}\!\left(
    \Homsh_{\Osh_X}\!\left(
      \underline{L}(A_1)\otimes\cdots\otimes \underline{L}(A_r),
      \underline{L}(B_1)\otimes\cdots\otimes \underline{L}(B_s)
    \right)\right)}
\\[3pt]
\symup{T_{r}^{s*}(A_1,\ldots,A_r;B_1,\ldots,B_s)}
&
\symup{={}\underline{\underline{R}}^{*}\Gamma\!\left(
    \Homsh_{\Osh_X}\!\left(
      \underline{L}(A_1)\otimes\cdots\otimes \underline{L}(A_r),
      \underline{L}(B_1)\otimes\cdots\otimes \underline{L}(B_s)
    \right)\right)}
\end{aligned}
\right.
$}
\end{equation}
Dans cette formule, $\symup{\underline{L}(F)}$ désigne une résolution finie localement libre du
faisceau algébrique cohérent $\symup{F}$, et $\symup{\underline{\underline{R}}^{*}\Gamma(K)}$ désigne
l'hypercohomologie de l'espace $\symup{X}$ par rapport au complexe de faisceaux $\symup{K}$. Si
$\symup{r}$ ou $\symup{s}$ est nul, on remplace le produit tensoriel des $\symup{\underline{L}(A_i)}$, resp. des
$\symup{\underline{L}(B_j)}$, par $\symup{\Osh_X}$. En particulier $\symup{\underline{T}_{0}^{0}}$ et $\symup{T_{0}^{0}}$
sont des foncteurs gradués à 0 arguments, $\symup{\underline{T}_{0}^{0}}$ se réduit au
degré 0, où il est le faisceau $\symup{\Osh_X}$, tandis que $\symup{T_{0}^{0}}$ est identique à
$\symup{H^{*}(X,\Osh_X)}$. Le fait que les deuxièmes membres de (8.1) ne dépendent pas
des résolutions choisies est de toutes façons évident pour la première ligne
(la question étant alors locale), et pour la deuxième résulte des remarques
générales qui précédaient, compte tenu de la suite spectrale
d'hypercohomologie du complexe de faisceaux
$\symup{K=\Homsh_{\Osh_X}\bigl(\underline{L}(A_1)\otimes\cdots,
      \underline{L}(B_1)\otimes\cdots\bigr)}$, aboutissant à l'hypercohomologie de $\symup{X}$ par rapport à $\symup{K}$, et dont le terme
initial est $\symup{H^{p}(X,H^{q}(K))}$, i.e.
\begin{equation}\tag{8.2}
  \symup{E_{2}^{p,q}
  =H^{p}\!\left(X,
    (\underline{T}_{r}^{s})^{(q)}
      (A_1,\ldots,A_r;B_1,\ldots,B_s)\right)}
\end{equation}
On voit alors que cette suite spectrale elle-même ne dépend pas des résolutions
choisies, son aboutissement est le premier membre de (8.1). On définit
facilement les opérateurs cobords relatifs aux divers arguments $\symup{A_i}$, $\symup{B_j}$,
en notant que toute suite exacte
$\symup{0\longrightarrow F'\longrightarrow F\longrightarrow F''\longrightarrow0}$
peut se résoudre par une suite exacte de complexes finis localement libres.

On définit dans le système des foncteurs $\symup{\underline{T}_{r}^{s*}}$, resp.
$\symup{T_{r}^{s*}}$, des opérations analogues à celles du calcul tensoriel, et dont la
définition est immédiate à partir des formules de définition (8.1). Ainsi, on a
une \uline{composition} (généralisant celle envisagée dans le n\textsuperscript{o} 2) :
\begin{equation}\tag{8.3}
\resizebox{\textwidth}{!}{$\symup{
T_{r}^{s*}(A_1,\ldots,A_r;B_1,\ldots,B_s)
\times T_{r'}^{s'*}(A'_1,\ldots,A'_{r'};B'_1,\ldots,B'_{s'})
\longrightarrow T_{r+r'}^{s+s'*}
(A_1,\ldots,A'_{r'};B_1,\ldots,B'_{s'})
}$}
\end{equation}
satisfaisant aux propriétés évidentes d'associativité, de compatibilité avec
les homomorphismes fonctoriels et les homomorphismes cobords, les suites
spectrales. De même, on a des opérations de symétrie qu'on laisse au lecteur le
soin d'expliciter. On a de plus une opération de \uline{contraction} chaque fois que
l'un des arguments $\symup{A_i}$

\src{191}
est égal à l'un des arguments $\symup{B_j}$ ;
\begin{equation}\tag{8.4}
\resizebox{\textwidth}{!}{$
\begin{aligned}
&
\symup{T_{r}^{s*}(A_1,\ldots,A_{i-1},C,A_{i+1},\ldots,A_r;
 B_1,\ldots,B_{j-1},C,B_{j+1},\ldots,B_s)\longrightarrow}
\\
&
\symup{\qquad\relbar\quad
T_{r-1}^{s-1*}(A_1,\ldots,\widehat{A}_i,\ldots,A_r;
 B_1,\ldots,\widehat{B}_j\ldots,B_s)}
\end{aligned}
$}
\end{equation}
De plus, si un argument $\symup{A_i}$ est un faisceau localement libre, alors on peut
l'enlever à condition de remplacer l'un des $\symup{A_j}$ ($\symup{j\neq i}$) par
$\symup{A_j\otimes A_i}$, ou l'un des $\symup{B_k}$ par $\symup{B_k\otimes A_i'}$ (où
$\symup{A_i'=\Homsh_{\Osh_X}(A_i,\Osh_X)}$), et on a une règle de calcul analogue dans
le cas où l'un des arguments $\symup{B_j}$ est localement libre. En particulier, on
peut toujours enlever un argument qui est identique à $\symup{\Osh_X}$. Si tous les
arguments sont localement libres, sauf au plus un des arguments $\symup{B_i}$, la règle
qu'on vient d'énoncer donne un isomorphisme fonctoriel
\begin{equation}\tag{8.5}
  \symup{T_{r}^{s*}(A_1,\ldots,A_r;B_1,\ldots,B_s)
  =H^{*}\bigl(X,A_1'\otimes\cdots\otimes A_r'
    \otimes B_1\otimes\cdots\otimes B_s\bigr)}
\end{equation}
(car on est ramené au cas où $\symup{r=0}$, $\symup{s=1}$, où cela est immédiat : on peut
aussi utiliser directement la suite spectrale de terme initial (8.2)). Des
opérations correspondant aux précédentes se définissent pour les
$\symup{\underline{T}_{r}^{s}}$. Les relations entre les diverses opérations ainsi
introduites sont les mêmes que pour les opérations analogues en calcul
tensoriel.

Soit $\symup{n}$ la dimension de $\symup{X}$. En appliquant successivement une composition
tensorielle (8.3) et des contractions (8.4) sur les arguments répétés, on
obtient un accouplement
\begin{equation}\tag{8.6}
\resizebox{\textwidth}{!}{$\symup{
  (T_{r}^{s})^{p}(A_1,\ldots;B_1,\ldots)
  \times
  (T_{s}^{r})^{n-p}(B_1,\ldots;A_1,\ldots,A_r\otimes\underline{\Omega}_X^{n})
  \longrightarrow H^{n}(X,\underline{\Omega}_X^{n})
}$}
\end{equation}

\result{THÉORÈME}{6} \uline{Si} $\symup{X}$ \uline{est une variété projective non singulière, alors
les accouplements} (8.6) \uline{sont des dualités}.

Cela résulte de façon purement formelle du corollaire au théorème 3. Il résulte
en effet facilement de ce corollaire que si $\symup{K}$ est un complexe de faisceaux
algébriques cohérents \uline{localement libres}, alors l'hypercohomologie de $\symup{X}$ par
rapport à $\symup{K}$ est en dualité avec l'hypercohomologie de $\symup{X}$ par rapport à
$\symup{K'\otimes\underline{\Omega}_X^{n}}$ par les accouplements naturels
\begin{equation}\tag{8.7}
  \symup{\underline{\underline{R}}^{p}\Gamma(K)
  \times\underline{\underline{R}}^{n-p}\Gamma(K'\otimes\underline{\Omega}_X^{n})
  \longrightarrow\underline{\underline{R}}^{n}\Gamma(\underline{\Omega}_X^{n})
  =H^{n}(X,\underline{\Omega}_X^{n})}
\end{equation}
On le voit en utilisant la suite spectrale de terme initial
$\symup{H^{p}(H^{q}(X,K))}$ et la suite spectrale analogue pour
$\symup{K'\otimes\underline{\Omega}_X^{n}}$. Du résultat précédent, le théorème 6 se déduit en
utilisant la définition (8.1).

\src{192}
\subsection*{\normalfont REMARQUES.}

\noindent 1\textsuperscript{o} Pour les définitions précédant le théorème 6,
il n'était pas nécessaire que $\symup{X}$ soit non singulière, car il n'était pas
indispensable de travailler seulement avec des résolutions \uline{finies}. Mais si $\symup{X}$
est singulière on ne pourra plus affirmer, a priori, que les
$\symup{(\underline{T}_{r}^{s})^{p}(A_1,\ldots;B_1,\ldots)}$ soient des faisceaux
\uline{cohérents}, car dans le complexe de faisceaux
\[
  \symup{\Homsh_{\Osh_X}\bigl(\underline{L}(A_1)\otimes\cdots,\underline{L}(B_1)\otimes\cdots\bigr)}
\]
il y aura une infinité de composantes de degré total donné.

\medskip
\noindent 2\textsuperscript{o} On vérifie facilement que dans les formules
(8.1), on peut remplacer \uline{un} des $\symup{\underline{L}(B_i)}$ par $\symup{B_i}$. Compte tenu de la
proposition 3, cela montre donc qu'on a
\begin{equation}\tag{8.8}
\left\{
\begin{aligned}\symup{\underline{T}_{1}^{1*}(A;B)}&\symup{=\Extsh_{\Osh_X}^{*}(A,B)}\\\symup{T_{1}^{1*}(A;B)}&\symup{=\operatorname{Ext}_{\Osh_X}^{*}(X;A,B)}\end{aligned}
\right.
\end{equation}
En particulier faisant dans (8.6) $\symup{r=s=1}$ et $\symup{A_1=\Osh_X}$, on retrouve le
théorème 3. La formule (8.8) implique aussi
$\symup{T_{0}^{1*}(B)=H^{*}(X,B)}$, $\symup{T_{1}^{0*}(A)=\operatorname{Ext}_{\Osh_X}^{*}(X;A,\Osh_X)}$

\medskip
\noindent 3\textsuperscript{o} On voit sur la formule (8.1) qu'en général les
foncteurs $\symup{\underline{T}_{r}^{s*}}$ et $\symup{T_{r}^{s*}}$ ont des composantes de
degrés positifs \uline{et} négatifs. Utilisant la remarque 2, on voit que si la
dimension de $\symup{X}$ est $\symup{n}$, alors les composantes non nulles de
$\symup{\underline{T}_{r}^{s*}}$ sont comprises entre $\symup{-(s-1)n}$ et $\symup{rn}$, si $\symup{s>0}$,
entre $\symup{0}$ et $\symup{rn}$, si $\symup{s=0}$, et les composantes non nulles de $\symup{T_{r}^{s*}}$
sont comprises entre $\symup{-(s-1)n}$ et $\symup{(r+1)n}$, si $\symup{s>0}$, entre $\symup{0}$ et $\symup{(r+1)n}$,
si $\symup{s=0}$, (et sauf erreur même, si $\symup{r>0}$, entre $\symup{-(s-1)n}$ et $\symup{rn}$, resp. entre
$\symup{0}$ et $\symup{rn}$).

\section*{\centering\normalfont BIBLIOGRAPHIE}

\noindent[1] CARTIER (Pierre). - Les groupes
$\symup{\operatorname{Ext}^{s}(A,B)}$, Séminaire A. Grothendieck : Algèbre homologique,
t. 1, 1957, n\textsuperscript{o} 3.

\noindent[2] GROTHENDIECK (Alexandre). - Sur quelques points d'algèbre
homologique, Tohoku math. J., t. 9, 1957, p. 119-183.

\noindent[3] SERRE (Jean-Pierre). - Faisceaux algébriques cohérents,
Annals of Math., t. 61, 1955, p. 197-278.

\noindent[4] SERRE (Jean-Pierre). - Sur la dimension homologique des anneaux
et des modules noethériens, Proc. Intern. Symp. on alg. number Theory
[1955. Tokyo et Nikko] - Tokyo, Science Council of Japan, 1956 ;
p. 175-189.

\clearpage
\src{193}
\section*{\centering\normalfont ADDITIF}

Les difficultés signalées dans la remarque de la page 13 sont actuellement
complètement résolues, grâce à une extension des théorèmes de dualité à des
variétés quelconques (à singularités arbitraires). Pour cette théorie, qui peut
se formuler dans le cadre général de la \fquote{théorie des schémas}, nous renvoyons
aux \fquote{Eléments de Géométrie algébrique} de J. DIEUDONNÉ et A. GROTHENDIECK,
actuellement en préparation. On trouvera quelques indications dans :

\begin{quote}
GROTHENDIECK (A.). - Cohomology theory of algebraic varieties,
Congrès int. des Mathématiciens, 1958, Edinburgh (A paraître).
\end{quote}

\begin{flushright}
[Avril 1959]
\end{flushright}

\begin{center}\rule{36mm}{0.5pt}\end{center}

\par\smallskip
\begingroup\footnotesize
\endgroup
